\documentclass[aspectratio=169,11pt,t]{beamer}
\input{../../../../theme/ksa-cs-v2.tex}
% 원본 덱 (ksa-theme): slides/v2/ml1/week01/data_handling/ -- 그림은 거기 figs/ 를 그대로 쓴다
\setbeamerfont{author}{size=\normalsize}
\setbeamerfont{institute}{size=\small}
\graphicspath{{../../../../v2/ml1/week01/data_handling/figs/}{../../../../assets/}}
\makeatletter
\newcommand{\markpage}[1]{\protected@write\@auxout{}{\string\newlabel{#1}{{\the\c@framenumber}{}}}\hypertarget{#1}{}}
\newcommand{\pref}[1]{\hyperlink{#1}{\textcolor{ksagolddk}{(p.\ref{#1})}}}
\makeatother

\renewcommand{\ksacredit}{Machine Learning and Deep Learning \,\textperiodcentered\, 데이터 직접 다루기}
\title{데이터 직접 다루기}
\subtitle{더티 데이터 (결측치 · 중복 · 이상치) 와 정규화 · 표준화}
\author{Machine Learning and Deep Learning}
\institute{한국과학영재학교 (KSA)}
\date{}
\begin{document}
{\setbeamertemplate{footline}{}\begin{frame}[c,noframenumbering]\titlepage\end{frame}}

\section{질문}
% --- 오늘의 질문
\begin{frame}{오늘의 질문}
\vfill
\centering\includegraphics[width=0.74\linewidth]{dd_hook.png}\par\vskip0.3em
{\large 우리 반 10명의 키. \kq{평균 키가 131cm?} 뭔가 이상하다}
\vfill
\note{우리 반 10명의 키를 모아서 평균을 냈더니 131센티미터가 나왔습니다. 고등학생 평균 키로는 너무 작죠. 그래프를 보면 이유가 보입니다. 수아와 준호 칸이 0입니다. 키를 안 적은 두 칸에 0이 들어가서 평균을 끌어내린 겁니다. 실제 데이터에는 이런 지저분한 값이 늘 섞여 있습니다. 이걸 더티 데이터라고 합니다. 오늘은 모델을 만들기 전에 데이터를 직접 열어 봅니다. 앞부분은 더티 데이터의 세 가지 모양, 뒷부분은 값의 크기를 맞추는 정규화입니다.}
\end{frame}

\section{더티 데이터}
% --- 세 가지 모양
\begin{frame}{더티 데이터의 세 가지 모양}
\vfill
\centering
{\Large
\begin{tabular}{ccc}
\kb{① 결측치} & \kb{② 중복} & \kb{③ 이상치}\\[0.4em]
\normalsize 있어야 할 값이 없다 & \normalsize 같은 기록이 여러 번 & \normalsize 혼자 크게 튀는 값\\[0.2em]
\normalsize Missing value & \normalsize Duplicate & \normalsize Outlier
\end{tabular}}\vskip1.4em
{\large 모델을 고르기 전에, 데이터부터 본다}\\[0.4em]
{\small Garbage in, garbage out: 쓰레기를 넣으면 쓰레기가 나온다}
\vfill
\note{더티 데이터는 모양이 딱 세 가지입니다. 있어야 할 값이 없는 결측치, 같은 기록이 여러 번 들어간 중복, 혼자 크게 튀는 이상치. 아무리 좋은 모델을 써도 데이터가 더러우면 결과도 더럽습니다. 영어로 garbage in, garbage out이라고 합니다. 그래서 모델을 고르기 전에 데이터부터 봐야 합니다.}
\end{frame}

% --- 결측치
\begin{frame}{① 결측치: 0을 빼면}
\vfill
\centering\includegraphics[width=0.74\linewidth]{dd_clean.png}\par\vskip0.3em
{\large 키를 안 적은 두 칸을 빼고 계산하면 \kb{평균 164cm}}
\vfill
\note{첫째, 결측치입니다. 수아와 준호는 키를 안 적었는데 0이 들어갔습니다. 0은 키가 아니라 비었다는 표시입니다. 이 두 칸을 빼고 8명으로 계산하면 평균 164센티미터. 훨씬 그럴듯합니다.}
\end{frame}

% --- 숫자인 척하는 결측치
\begin{frame}{숫자인 척하는 결측치가 더 위험하다}
\markpage{sentinel}
\vfill
\centering
{\Large
\begin{tabular}{lll}
\textbf{적힌 값} & \textbf{컴퓨터가 보기에} & \\[0.3em]\hline\\[-0.5em]
(빈칸) & 비었음 & \kb{눈에 띈다}\\[0.4em]
\texttt{NA}, \texttt{?} & 글자 & 계산하다 오류가 난다\\[0.4em]
\texttt{0}, \texttt{-999} & \kq{진짜 숫자} & \kq{오류 없이 평균에 섞인다}
\end{tabular}}\vskip1.0em
{\small 키 0cm, 나이 $-999$세처럼 \textbf{말이 안 되는 값}을 찾아서 빈칸으로 바꾼다}
\vfill
\note{빈칸은 눈에 띄니까 오히려 괜찮습니다. 문제는 비었다는 뜻을 숫자로 적은 경우입니다. 0이나 마이너스 999는 컴퓨터가 보기에 진짜 숫자라서 오류도 안 나고 평균에 조용히 섞입니다. 그래서 키 0센티, 나이 마이너스 999세처럼 말이 안 되는 값을 직접 찾아서 진짜 빈칸으로 바꿔야 합니다.}
\end{frame}

% --- 중복
\begin{frame}{② 중복: 같은 기록이 여러 번}
\vfill
\centering\includegraphics[width=0.76\linewidth]{dd_dup.png}\par\vskip0.3em
{\large 민수 기록이 세 번 들어가면 \kq{평균 166cm} (원래 164cm)}
\vfill
\note{둘째, 중복입니다. 설문을 여러 번 제출했거나 파일을 합치다가 같은 기록이 여러 번 들어가는 경우입니다. 키가 큰 민수 기록이 세 번 들어가면 평균이 164에서 166으로 올라갑니다. 한 사람이 세 표를 행사한 셈입니다.}
\end{frame}

% --- 이상치
\begin{frame}{③ 이상치: 혼자 크게 튀는 값}
\vfill
\centering\includegraphics[width=0.76\linewidth]{dd_outlier.png}\par\vskip0.3em
{\large 준호 250cm는 \kq{오타}, 태오 205cm는 \kb{농구 선수}일 수도}
\vfill
\note{셋째, 이상치입니다. 준호의 키가 250센티미터로 적혀 있습니다. 세계 최장신 기록에 가까운 값이니 오타일 가능성이 큽니다. 이런 오타는 결측치처럼 비워 둡니다. 그런데 태오의 205센티미터는 어떨까요? 농구 선수라면 진짜일 수 있습니다. 같은 튀는 값이라도 처리 방법이 다릅니다.}
\end{frame}

% --- 그림으로 찾기
\begin{frame}{그림으로 찾는다: 상자 그림}
\vfill
\centering\includegraphics[height=0.66\textheight]{dd_box.png}\par\vskip0.3em
{\large \kq{수염 밖에 홀로 떨어진 점}이 이상치 후보} \quad {\small 후보를 찾는 건 그림이, 지울지 남길지는 사람이 정한다}
\vfill
\note{이상치는 숫자 표만 봐서는 잘 안 보이고, 그림을 그리면 바로 보입니다. 왼쪽이 상자 그림 읽는 법입니다. 값을 작은 것부터 줄 세웠을 때 25퍼센트 지점이 상자의 아래 끝, 75퍼센트 지점이 위 끝이라서 상자 안에 가운데 절반이 들어 있습니다. 상자 속 빨간 선이 한가운데 값, 중앙값, 영어로 median입니다. 수염은 상자에서 상자 길이의 1.5배 안에 있는 가장 먼 값까지 그리고, 그보다 바깥에 있는 점을 따로 찍어 이상치 후보로 보여 줍니다. 오른쪽 우리 반 키를 보면 상자 길이가 12.5센티라서 기준선이 191.75이고, 205와 250이 그 위로 튀어나와 있습니다. 다만 그림은 후보를 알려 줄 뿐이고, 지울지 남길지는 사람이 정합니다.}
\end{frame}

% --- 지울까 남길까
\begin{frame}{지울까, 고칠까, 남길까}
\vfill
\centering
{\large
\begin{tabular}{lll}
\textbf{경우} & \textbf{예} & \textbf{처리}\\[0.3em]\hline\\[-0.5em]
분명한 오타 & 키 250cm & 빈칸으로 (결측치처럼)\\[0.4em]
단위 실수 & 꽃받침 51cm (mm 로 적음) & 고친다 ($\div 10$)\\[0.4em]
진짜일 수 있음 & 농구 선수 205cm & \kb{남긴다}\\[0.4em]
모르겠음 & ? & 표시해 두고, 지운 결과와 둘 다 보고
\end{tabular}}\vskip1.0em
\kq{판단의 근거는 그 분야를 아는 지식 (도메인 지식)}
\vfill
\note{튀는 값을 어떻게 할지 정리하면 이렇습니다. 분명한 오타는 결측치처럼 비웁니다. 단위를 잘못 적은 경우는 고칠 수 있습니다. 실습에서 보겠지만 꽃받침 길이 51은 센티미터가 아니라 밀리미터로 적은 겁니다. 진짜일 수 있는 값은 남깁니다. 정말 모르겠으면 표시해 두고 지운 결과와 안 지운 결과를 둘 다 보고합니다. 이 판단의 근거가 그 분야를 아는 지식, 도메인 지식입니다.}
\end{frame}

% --- 모델도 속는다
\begin{frame}{모델도 속는다}
\vfill
\centering\includegraphics[width=0.66\linewidth]{dd_iris.png}\par\vskip0.3em
{\large 붓꽃 150송이 중 \kq{6송이}의 꽃잎 너비를 0으로 적었을 뿐인데, 직선이 끌려 내려간다}
\vfill
\note{더러운 데이터는 평균만 망치는 게 아니라 모델도 속입니다. 붓꽃 150송이의 꽃잎 길이와 너비입니다. 그중 6송이의 너비를 비었다는 뜻으로 0이라고 적었습니다. 초록 선은 깨끗한 데이터로 그은 직선, 빨간 선은 0이 섞인 데이터로 그은 직선입니다. 150송이 중 6송이일 뿐인데 기울기가 0.42에서 0.37로 내려갔습니다. 모델은 0을 진짜 너비로 믿고 배웁니다.}
\end{frame}

\section{크기 맞추기}
% --- 정규화 질문
\begin{frame}{두 번째: 값의 크기를 맞춘다}
\vfill
\centering\includegraphics[width=0.6\linewidth]{fs_hook.png}\par\vskip0.3em
{\large 같은 신경망, 같은 와인 데이터. \kq{성분 값을 그대로 넣으면 72\%}, 정규화하면 \kb{98\%}}
\vfill
\note{와인 178병의 성분을 보고 품종 세 가지를 맞히는 신경망입니다. 모델도 같고 학습 방법도 같습니다. 바뀐 건 하나, 입력 값을 그대로 넣었느냐, 정규화해서 넣었느냐입니다. 그대로 넣으면 시험 정확도가 72퍼센트, 정규화하면 98퍼센트입니다. 시드를 10번 바꿔 평균 낸 값입니다. 숫자 몇 개를 빼고 나눈 것뿐인데 왜 이렇게 차이가 날까요? 이제 그 이유를 보고, 이미지, 글, 시계열, 소리, 비디오 데이터에서는 각각 무엇을 맞춰야 하는지 직접 데이터를 열어 봅니다.}
\end{frame}

% --- 성분 범위
\begin{frame}{와인 한 병 = 성분 값 13가지, 범위가 제각각}
\vfill
\centering\includegraphics[width=0.74\linewidth]{fs_ranges.png}\par\vskip0.2em
{\large \kq{프롤린 278 $\sim$ 1,680}, \textcolor{ksagolddk}{\textbf{색조 0.48 $\sim$ 1.71}}: 천 배 넘게 차이 난다}
\vfill
\note{와인 한 병마다 알코올, 마그네슘, 프롤린, 색조 같은 성분 값이 13가지 있습니다. 가로축이 로그 눈금이라 한 칸이 10배입니다. 맨 위 빨간 막대 프롤린은 278에서 1,680, 아래쪽 색조는 0.48에서 1.71입니다. 단위가 다르니 크기도 제각각입니다. 사람은 단위를 보고 알아서 읽지만, 모델은 숫자만 봅니다.}
\end{frame}

% --- 뉴런
\begin{frame}{뉴런 하나가 하는 일}
\vfill
\centering
\begin{columns}[c]
\column{0.42\textwidth}\centering
\includegraphics[width=\linewidth]{fs_neuron.png}
\column{0.54\textwidth}
{\LARGE\[ \text{output} = \tanh(\textcolor{ksablue}{w}\cdot x + b) \]}\vskip0.6em
{\large $x$ = 성분 값 13개 (데이터)\\[0.4em]
$\textcolor{ksablue}{w}$ = 뉴런이 \kb{학습할 가중치} 13개}\\[0.8em]
PyTorch \texttt{nn.Linear(13, 16)} 의 $w$ 는 처음에\\
\kb{$-0.28 \sim 0.28$} 사이 작은 난수 ($\pm 1/\sqrt{13}$)
\end{columns}
\vfill
\note{뉴런은 뇌세포처럼 입력 x에 가중치 w를 곱해 다 더하고, 탄젠트 에이치 같은 곡선을 거쳐 내보냅니다. 여기서 x는 데이터, 와인 성분 값이고, w는 뉴런이 학습할 변수입니다. 파이토치의 nn.Linear는 w를 마이너스 0.28에서 0.28 사이 작은 난수로 시작합니다. 입력이 13개라서 1 나누기 루트 13입니다. 즉 처음 w는 0 근처의 작은 수인데, x는 1,000이 넘을 수 있습니다.}
\end{frame}

% --- 계산
\begin{frame}{와인 1번 병을 그대로 넣으면}
\markpage{rawsum}
\vfill
\centering
{\large
\begin{tabular}{lrcrcr}
\textbf{성분} & \textbf{값 $x$} & & \textbf{가중치 $w$} & & \textbf{$x \times w$}\\\hline\\[-0.8em]
프롤린 & \kq{1,065} & $\times$ & 0.15 & $=$ & \kq{162.15}\\
마그네슘 & 127 & $\times$ & 0.05 & $=$ & 6.56\\
알코올 & 14.23 & $\times$ & 0.11 & $=$ & 1.57\\
색조 & 1.04 & $\times$ & $-0.16$ & $=$ & $-0.17$\\
나머지 9개 $+\,b$ & & & & & $-1.98$\\\hline\\[-0.8em]
\multicolumn{5}{l}{\textbf{합} $w\cdot x + b$} & \kq{168.13}
\end{tabular}}\vskip0.8em
{\Large $\tanh(168.13) = \kq{1.00}$}\quad 프롤린 하나가 합을 다 차지한다\\[0.3em]
{\scriptsize $w$ 는 반올림해 적었고, 곱은 반올림 전 값으로 계산 (seed 0, 4번 뉴런)}
\vfill
\note{와인 1번 병을 그대로 넣어 봅시다. 프롤린 1,065에 가중치 0.15를 곱하면 162입니다. 마그네슘은 6.6, 알코올은 1.6, 색조는 0.2도 안 됩니다. 나머지 아홉 개와 b를 다 더해도 마이너스 2 정도. 합은 168입니다. 탄젠트 에이치에 168을 넣으면 1.00, 곡선 끝에 붙어 버립니다. 사실상 프롤린 하나만 보고 있는 셈입니다.}
\end{frame}

% --- 포화
\begin{frame}{12병 모두 출력 1.00: 기울기 0, 배우지 못한다}
\vfill
\centering\includegraphics[width=0.74\linewidth]{fs_tanh_raw.png}\par\vskip0.3em
{\large 곡선이 평평한 곳 = \kq{기울기 0}. 학습 전 뉴런 출력의 \kq{98\%}가 이렇게 막혀 있다}
\vfill
\note{1번 병만 그런 게 아닙니다. 다른 와인 12병을 넣어 보면 합이 45에서 200 사이, 출력은 모두 1.00입니다. 곡선이 평평한 곳에 있으니 기울기가 0입니다. 기울기는 w를 어느 쪽으로 얼마나 고칠지 알려 주는 신호인데, 이게 0이면 학습이 안 됩니다. 학습 시작 때 은닉층 출력의 98퍼센트가 절댓값 0.99를 넘는, 이렇게 막힌 상태였습니다.}
\end{frame}

% --- 맞추면
\begin{frame}{성분마다 분포를 맞추면}
\vfill
\centering\includegraphics[width=0.74\linewidth]{fs_tanh_z.png}\par\vskip0.3em
{\large 프롤린 1,065 $\to$ \kb{1.01}. 합 \kb{0.67}, 출력 \kb{0.58}: 곡선의 가파른 곳에서 신호가 산다}
\vfill
\note{이번에는 성분마다 분포를 맞춰서 넣습니다. 방법은 곧 보겠습니다. 프롤린 1,065는 1.01이 되고, 다른 성분도 모두 0 근처의 비슷한 크기가 됩니다. 1번 병의 합은 0.67, 출력은 0.58입니다. 12병이 곡선의 가파른 곳에 고루 퍼집니다. 기울기가 살아 있으니 학습이 됩니다.}
\end{frame}

% --- Min-Max
\begin{frame}{방법 1. Min-Max 정규화: 0 $\sim$ 1 로}
\vfill
\centering
{\LARGE\[ x' = \frac{x - \textcolor{ksablue}{\min}}{\textcolor{ksablue}{\max} - \textcolor{ksablue}{\min}} \]}\vskip0.4em
{\large 성분마다 따로: 가장 작은 값 $\to$ 0, 가장 큰 값 $\to$ 1}\vskip1em
{\Large 프롤린 1,065: \quad $\dfrac{1065 - \textcolor{ksablue}{278}}{\textcolor{ksablue}{1680} - \textcolor{ksablue}{278}} = \dfrac{787}{1402} = \kb{0.56}$}\vskip1em
{\small 쓰기 좋은 곳: 이미지 픽셀처럼 \textbf{범위가 정해진 값} (0 $\sim$ 255 $\to$ 0 $\sim$ 1)}
\vfill
\note{첫 번째 방법은 Min-Max 정규화입니다. 성분마다 따로, 가장 작은 값을 빼고 범위로 나눕니다. 그러면 가장 작은 값은 0, 가장 큰 값은 1이 됩니다. 프롤린 1,065는 278을 빼서 787, 범위 1,402로 나누면 0.56입니다. 이미지 픽셀처럼 0에서 255로 범위가 정해진 값에 잘 맞습니다.}
\end{frame}

% --- Z-score
\begin{frame}{방법 2. Z-score 표준화: 평균 0, 표준편차 1 로}
\markpage{zscore}
\vfill
\centering
{\LARGE\[ z = \frac{x - \textcolor{ksablue}{\mu}}{\textcolor{ksared}{\sigma}} \]}\vskip0.2em
{\large $\textcolor{ksablue}{\mu}$ = 평균 (mean), \quad $\textcolor{ksared}{\sigma}$ = 표준편차 (standard deviation), \quad 성분마다 따로}\vskip1em
{\Large 프롤린 1,065: \quad $\dfrac{1065 - \textcolor{ksablue}{747}}{\textcolor{ksared}{314}} = \kb{1.01}$}\vskip0.6em
{\large ``평균보다 \textbf{표준편차 1개만큼} 많다''}\vskip0.8em
{\small 보통은 이것을 쓴다. 튀는 값 하나가 있어도 Min-Max 보다 덜 흔들린다}
\vfill
\note{두 번째는 Z-score 표준화입니다. 성분마다 평균을 빼고 표준편차로 나눕니다. 프롤린의 178병 평균은 747, 표준편차는 314입니다. 1,065에서 747을 빼면 318, 314로 나누면 1.01. 평균보다 표준편차 한 개만큼 많다는 뜻입니다. 결과는 평균 0, 표준편차 1이 됩니다. Min-Max는 가장 큰 값 하나가 튀면 나머지가 전부 0 근처로 눌리는데, Z-score는 덜 흔들립니다. 그래서 특별한 이유가 없으면 보통 Z-score를 씁니다. 참고로 여기 표준편차는 n으로 나눈 값, numpy의 std 기본값입니다.}
\end{frame}

% --- 학습 데이터로만
\begin{frame}{평균 · 표준편차는 학습 데이터로만 구한다}
\vfill
\centering
{\large
\begin{tabular}{lll}
\textbf{단계} & \textbf{할 일} & \\[0.2em]\hline\\[-0.6em]
학습 데이터 & $\mu, \sigma$ (또는 min, max) 를 \kb{구한다} & \texttt{fit}\\[0.4em]
학습 데이터 & 그 값으로 바꾼다 & \texttt{transform}\\[0.4em]
시험 · 새 데이터 & \kb{같은} $\mu, \sigma$ 로 바꾼다 (새로 구하지 않는다) & \texttt{transform}
\end{tabular}}\vskip1em
{\large 시험 데이터로 평균을 구하면 \kq{시험 문제를 미리 엿본 것} (데이터 누수)}\vskip0.8em
{\small 회귀에서 정답 $y$ 도 표준화했다면, 예측값은 마지막에 $\hat y = \hat z\,\sigma_y + \mu_y$ 로 원래 단위로 되돌린다}
\vfill
\note{중요한 규칙이 하나 있습니다. 평균과 표준편차는 학습 데이터로만 구합니다. 사이킷런으로 말하면 학습 데이터에 fit, 학습과 시험 데이터 모두에 같은 값으로 transform입니다. 시험 데이터로 평균을 다시 구하면 시험 문제를 미리 엿본 셈이 됩니다. 이걸 데이터 누수라고 합니다. 실제로 서비스할 때 새로 들어오는 데이터는 한 개씩 오니까 평균을 구할 수도 없습니다. 또, 회귀에서 정답 y까지 표준화했다면 모델이 내놓은 값은 표준화된 값이니, 마지막에 표준편차를 곱하고 평균을 더해 원래 단위로 되돌려야 합니다.}
\end{frame}

\section{데이터마다}
% --- 데이터마다
\begin{frame}{데이터마다: ``무엇마다'' 맞추나}
\markpage{overview}
\vfill
\centering
{\normalsize
\begin{tabular}{llll}
\textbf{데이터} & \textbf{모양} & \textbf{무엇마다} & \textbf{보통 하는 일}\\[0.2em]\hline\\[-0.6em]
표 (tabular) & (행, 열) & \kb{열(성분)마다} & Z-score\\[0.35em]
흑백 이미지 & (H, W) & 이미지 전체 & $\div 255$, 그다음 평균 · 표준편차\\[0.35em]
컬러 이미지 & (H, W, 3) & \kb{채널(R, G, B)마다} & $\div 255$, 채널별 평균 · 표준편차\\[0.35em]
언어 & 단어의 줄 & 글자 표기 & 표기 맞추기, 번호는 \kq{그대로}\\[0.35em]
시계열 & (시간,) & \kb{과거만}, 또는 창마다 & 학습 기간 Z-score, 창마다 Z-score\\[0.35em]
소리 & (샘플,) & 녹음마다 & $\div 32768$, 음량 맞추기, dB\\[0.35em]
비디오 & (T, H, W, 3) & 채널마다, \kb{영상 전체로} & 이미지와 같게, 프레임마다 따로는 X
\end{tabular}}\vskip0.8em
{\large 먼저 \kq{데이터를 열어서} 모양 · 자료형 · 최솟값 · 최댓값을 본다}
\vfill
\note{지금까지는 표 데이터였습니다. 열마다 단위가 다르니 열마다 맞췄습니다. 다른 데이터는 어떨까요? 핵심 질문은 무엇마다 맞추느냐입니다. 흑백 이미지는 모든 픽셀이 같은 단위라 이미지 전체를 한꺼번에, 컬러는 R, G, B 채널마다, 언어는 숫자 크기가 아니라 표기를 맞추고, 시계열은 과거만 보고, 소리는 녹음마다 음량을 맞추고, 비디오는 영상 전체로 맞춥니다. 하나씩 실제 데이터를 열어 보겠습니다. 어떤 데이터든 첫 단계는 같습니다. 모양, 자료형, 최솟값, 최댓값을 직접 찍어 보는 겁니다.}
\end{frame}

% --- 흑백 이미지
\begin{frame}{흑백 이미지: 픽셀 0 $\sim$ 255 를 그대로 넣으면}
\vfill
\centering\includegraphics[width=0.82\linewidth]{fs_mnist.png}\par\vskip0.2em
{\large MNIST: $\div 255$ 만 해도 \kq{63\%} $\to$ \kb{90\%}. 그대로면 학습 전 뉴런 \kq{95\%}가 막힌다}\\[0.3em]
{\scriptsize 학습 1만 장 · 시험 1만 장, MLP 784-64-10 (tanh), SGD 0.1, 3에폭, 시드 3개 평균. 0.1307, 0.3081 = 학습 6만 장의 픽셀 평균 · 표준편차}
\vfill
\note{손글씨 숫자 MNIST입니다. 28 곱하기 28 픽셀, 값은 0에서 255 사이의 정수, 자료형은 uint8입니다. 모든 픽셀이 같은 단위라서 이미지 전체를 같은 수로 나눕니다. 그대로 넣으면 와인처럼 뉴런 95퍼센트가 막히고 정확도 63퍼센트. 255로 나누기만 해도 90퍼센트, 학습 데이터의 평균 0.1307과 표준편차 0.3081로 표준화하면 92퍼센트입니다. 한 가지 함정: uint8끼리 빼면 음수가 안 돼서, 10 빼기 20이 246이 됩니다. 계산 전에 반드시 실수형으로 바꿔야 합니다. 실습에서 직접 확인합니다.}
\end{frame}

% --- 컬러 이미지
\begin{frame}{컬러 이미지: 채널(R, G, B)마다 따로}
\vfill
\centering\includegraphics[width=0.94\linewidth]{fs_cifar.png}\par\vskip0.3em
{\large CIFAR-10: 평균 R \textbf{0.494}, G \textbf{0.485}, B \textbf{0.450}. 채널마다 빼고 나눈다}\\[0.4em]
{\small 사진은 (H, W, C), PyTorch 는 \kb{(C, H, W)}: 축 순서부터 확인\\[0.2em] 남이 학습시킨 모델은 그 모델의 값을 그대로 (ImageNet: 0.485, 0.456, 0.406)}
\vfill
\note{컬러 이미지는 한 픽셀에 빨강, 초록, 파랑 세 값이 있습니다. CIFAR-10 시험 이미지 1만 장으로 채널별 평균을 내면 빨강 0.494, 초록 0.485, 파랑 0.450입니다. 하늘이나 바다처럼 파란 배경도 있지만 전체로는 파랑이 조금 어둡습니다. 그래서 채널마다 따로 평균을 빼고 표준편차로 나눕니다. 결과는 숫자 세 개씩입니다. 두 가지 주의할 점. 사진 파일을 열면 높이, 너비, 채널 순서인데 파이토치는 채널을 맨 앞에 둡니다. 그리고 남이 학습시킨 모델을 가져다 쓸 때는 그 모델이 학습할 때 쓴 값, 예를 들어 ImageNet의 0.485, 0.456, 0.406을 그대로 써야 합니다.}
\end{frame}

% --- 언어
\begin{frame}{언어: 숫자 크기가 아니라 \textbf{표기}를 맞춘다}
\vfill
\centering
{\normalsize
\begin{tabular}{ll}
\kb{① 표기 맞추기} & \texttt{FREE}, \texttt{Free}, \texttt{free!} $\to$ \texttt{free}\\
& SMS 5,574통의 단어 종류 \kq{15,733} $\to$ \kb{8,713}\\[0.8em]
\kb{② 단어 번호는 크기가 아니다} & \texttt{free} $= 3373$, \texttt{call} $= 1828$: 3373 이 ``더 큰'' 단어는 아니다\\
& $\to$ 번호를 표준화하지 않는다. 임베딩 벡터로 바꾼다\\
& (\texttt{nn.Embedding} 은 처음에 평균 0, 표준편차 1 난수)\\[0.8em]
\kb{③ 단어 세기는 길이로 나누기} & 1단어 문자와 171단어 문자: 긴 글일수록 숫자가 커진다
\end{tabular}}
\vfill
\note{언어 데이터는 다릅니다. 첫째, 숫자가 아니라 글자의 표기를 맞춥니다. 대문자 FREE, 첫 글자만 대문자 Free, 느낌표 붙은 free는 모두 같은 단어입니다. 영어 문자 메시지 5,574통을 공백으로만 자르면 단어 종류가 15,733개인데, 소문자로 바꾸고 문장부호를 떼면 8,713개로 줄어듭니다. 둘째, 단어마다 번호를 붙이지만 그 번호는 크기가 아닙니다. free가 3373번이고 call이 1828번이라고 free가 더 큰 단어가 아니죠. 그래서 번호는 표준화하지 않고, 임베딩이라는 벡터로 바꿉니다. 파이토치 임베딩은 처음에 평균 0, 표준편차 1인 난수로 시작합니다. 이미 표준화된 크기에서 출발하는 셈입니다. 셋째, 단어를 세어 벡터로 만들면 긴 글일수록 숫자가 커지니, 길이로 나눠 비율로 만듭니다.}
\end{frame}

% --- 시계열
\begin{frame}{시계열: 미래는 과거의 범위를 넘는다}
\vfill
\centering\includegraphics[width=0.8\linewidth]{fs_air.png}\par\vskip0.2em
{\large 학습 최대 505 $\to$ 시험 \kq{622} (Min-Max \kq{1.29}) \quad 창마다 표준화: 오차 \kq{74} $\to$ \kb{10}천 명}\\[0.3em]
{\scriptsize 국제선 항공 승객 1949$\sim$1960, 앞 10년 학습 · 뒤 2년 시험, 지난 12달 $\to$ 다음 달, $k$-최근접 이웃 ($k = 3$)}
\vfill
\note{시계열은 시간 순서가 있는 데이터입니다. 1949년부터 1960년까지 한 달 국제선 승객 수입니다. 앞 10년으로 학습하고 뒤 2년을 예측합니다. 규칙대로 학습 기간으로만 평균과 표준편차를 구합니다. 그런데 승객은 계속 늘어서, 학습 기간 최대가 505인데 시험 기간에는 622까지 갑니다. Min-Max로 바꾸면 1.29, 1을 넘습니다. 과거에서 가장 비슷한 패턴을 찾아오는 최근접 이웃 모델은 본 적 없는 크기를 내놓지 못해서, 빨간 점선처럼 낮게 예측합니다. 오차가 평균 74천 명. 해결책은 12달 창마다 그 창의 평균과 표준편차로 표준화하는 겁니다. 그러면 모델은 크기가 아니라 모양, 여름에 오르고 겨울에 내리는 패턴을 배웁니다. 예측한 뒤 그 창의 평균과 표준편차로 되돌리면 오차가 10천 명으로 줄어듭니다.}
\end{frame}

% --- 소리
\begin{frame}{소리: 녹음마다 음량이 다르다}
\vfill
\centering\includegraphics[width=0.84\linewidth]{fs_sound.png}\par\vskip0.2em
{\large 같은 ``zero'', 음량 \kq{37배} 차이. int16 $\div\, 32768 \to -1 \sim 1$, 그다음 \kb{녹음마다} 맞춘다}\\[0.3em]
{\scriptsize Free Spoken Digit Dataset, 8,000 Hz. 가장 큰 값으로 나누기 = 피크 정규화 (녹음마다 표준편차로 나누는 방법도 쓴다)}
\vfill
\note{소리 파일은 1초에 8,000번, 마이크가 받은 압력을 숫자로 적은 것입니다. 자료형은 16비트 정수 int16, 범위가 마이너스 32,768에서 32,767입니다. 먼저 32,768로 나눠 마이너스 1에서 1로 만듭니다. 그런데 같은 zero라는 말인데, 잭슨의 녹음은 가장 큰 값이 0.74, 테오는 0.02입니다. 37배 차이. 마이크와 거리 차이일 뿐 말의 내용은 같습니다. 그래서 녹음마다 가장 큰 값으로 나누면 둘 다 1이 되고 모양을 비교할 수 있습니다.}
\end{frame}

% --- 비디오
\begin{frame}{비디오: 이미지 $\times$ 시간, \textbf{영상 전체}로 맞춘다}
\vfill
\centering\includegraphics[width=0.76\linewidth]{fs_video.png}\par\vskip0.2em
{\large (T, H, W, C) = (15, 200, 300, 3). 프레임마다 표준화하면 \kq{어두워지는 변화가 사라진다}}\\[0.3em]
{\scriptsize Muybridge, \textit{The Horse in Motion} (1878). 실습에서는 일부러 점점 어둡게 만든 영상으로 비교}
\vfill
\note{비디오는 이미지를 시간 순서로 쌓은 것입니다. 1878년 머이브리지가 찍은 달리는 말, 15프레임, 높이 200, 너비 300, 채널 3. 처리는 이미지와 같이 255로 나누고 채널별 평균과 표준편차로 맞춥니다. 다만 그 평균을 어디서 구하느냐가 중요합니다. 첫 줄은 일부러 점점 어둡게 만든 영상입니다. 둘째 줄처럼 프레임마다 따로 표준화하면 모든 프레임이 같은 밝기가 되어 해가 지는 변화가 사라집니다. 셋째 줄처럼 영상 전체, 더 정확히는 학습 데이터 전체의 값으로 맞추면 변화가 남습니다. 밝기 변화가 중요한 문제라면 프레임마다 따로 맞추면 안 됩니다.}
\end{frame}

\section{연습}
% --- 문제풀이
\begin{frame}{문제풀이}
\vfill
\begin{enumerate}
\item 키 기록 \texttt{[160, 0, 170, 0, 180]}. 평균은? 0을 결측치로 보고 빼면?\vskip0.6em
\item 다음 값은 오타일까, 진짜일까? 무엇을 더 알아야 할까? \quad (a) 고등학생 몸무게 450kg \quad (b) 하루 공부 시간 25시간\vskip0.6em
\item 8명의 하루 공부 시간 \texttt{[2, 4, 4, 4, 5, 5, 7, 9]}. 평균과 표준편차($n$ 으로 나눔)는? 9시간의 Z-score 는? 5시간의 Min-Max 값은?\vskip0.6em
\item 학습 데이터의 키 최솟값 150, 최댓값 190 으로 Min-Max 를 만들었다. 시험 데이터에 키 200 이 나오면 값은? 시험 데이터로 min, max 를 다시 구해야 할까?\vskip0.6em
\item \texttt{np.uint8(10) - np.uint8(20)} 은 얼마일까? 컬러 이미지 \texttt{X.shape == (10000, 32, 32, 3)} 의 채널별 평균은 \texttt{X.mean(axis=?)} ?
\end{enumerate}
\vfill
\note{1번 답. 다섯 개 평균은 510 나누기 5, 102입니다. 0 두 개를 빼면 510 나누기 3, 170입니다. 2번 답. a는 45킬로그램에 0이 하나 더 붙은 오타일 가능성이 큽니다. b는 하루가 24시간이니 불가능한 값입니다. 3번 답. 합이 40이라 평균 5. 평균과의 차를 제곱하면 9, 1, 1, 1, 0, 0, 4, 16, 합 32, 8로 나누면 분산 4, 표준편차 2. 9시간은 9 빼기 5 나누기 2, Z-score 2입니다. 5시간의 Min-Max는 5 빼기 2 나누기 7, 7분의 3, 약 0.43. 4번 답. 200 빼기 150 나누기 40, 1.25입니다. 1을 넘어도 괜찮습니다. 시험 데이터로 다시 구하면 학습 때와 다른 자로 재는 셈이고 데이터 누수입니다. 5번 답. uint8은 0에서 255만 담을 수 있어서 246이 됩니다. 먼저 실수형으로 바꿉니다. 채널별 평균은 axis 0, 1, 2. 장, 높이, 너비 방향으로 평균을 내고 채널만 남기니 숫자 3개입니다.}
\end{frame}

% --- 실습
\begin{frame}[fragile]{실습: 데이터를 직접 열어 보기}
\vfill
\centering\begin{tabular}{rl}
\large 1 & 더티 데이터: 일부러 더럽힌 iris 청소 (빈칸 · 0 · $-999$, 표기, 51cm 꽃받침) \kq{86\%} $\to$ \kb{96\%}\\[0.35em]
\large 2 & 표: 와인 13성분, Min-Max · Z-score 구현 $\to$ \kq{72\%} vs \kb{98\%}\\[0.35em]
\large 3 & 흑백 이미지: MNIST uint8 함정, $\div 255$, 평균 · 표준편차\\[0.35em]
\large 4 & 컬러 이미지: CIFAR-10 채널별 평균 · 표준편차, (H, W, C) $\to$ (C, H, W)\\[0.35em]
\large 5 & 언어: SMS 표기 맞추기 (15,733 $\to$ 8,713), 단어 세기 길이로 나누기\\[0.35em]
\large 6 & 시계열: 항공 승객, 학습 기간으로만 vs 창마다 표준화 (\kq{74} $\to$ \kb{10})\\[0.35em]
\large 7 & 소리: 숫자 말소리 6명, $\div 32768$, 음량 맞추기, dB 스펙트로그램\\[0.35em]
\large 8 & 비디오: 달리는 말 GIF, 프레임마다 vs 영상 전체
\end{tabular}\vskip0.8em
데이터는 첫 셀에서 내려받기 (MNIST · CIFAR-10 은 수십 MB) \\[0.3em]\href{https://colab.research.google.com/github/SuminHan/book-ml/blob/main/slides/v2/ml1/week01/data_handling/lab/data_handling.ipynb}{\textcolor{ksagolddk}{▶ Colab에서 실습 노트북 열기}}
\vfill
\note{실습은 두 부분입니다. 먼저 제가 일부러 더럽혀 둔 붓꽃 데이터를 청소합니다. 빈칸, 0, 마이너스 999, 대소문자와 공백이 섞인 꽃 이름, 중복된 줄, 밀리미터로 적은 꽃받침 길이를 하나씩 찾아 고치면 정확도가 86에서 96퍼센트가 됩니다. 그다음 일곱 가지 데이터를 열어 정규화합니다. 순서는 매번 같습니다. 모양, 자료형, 최솟값, 최댓값을 찍어 보고, 무엇마다 맞출지 정한 뒤, 함수를 직접 만들어 바꾸고, 바꾼 뒤의 값을 다시 확인합니다.}
\end{frame}

\section{정리}
% --- 정리
\begin{frame}{정리}
\vfill
\centering\large\begin{tabular}{l}
결측치는 \texttt{0}, \texttt{-999} 처럼 \kq{숫자인 척}한다. 이상치는 그림으로 찾고 판단은 사람이\\[0.45em]
값의 크기가 제각각이면 뉴런이 \kq{막힌다} (tanh 출력 1.00, 기울기 0)\\[0.45em]
Min-Max $\dfrac{x-\min}{\max-\min}$ \quad Z-score $\dfrac{x-\mu}{\sigma}$ \quad 통계는 \kb{학습 데이터로만}\\[0.7em]
\kb{무엇마다} 맞출지는 데이터가 정한다: 열 · 채널 · 녹음 · 창 · 영상 전체\\[0.45em]
어떤 데이터든 먼저 \kq{열어서} 모양 · 자료형 · 최솟값 · 최댓값을 본다
\end{tabular}
\vfill
\note{정리합니다. 더티 데이터는 결측치, 중복, 이상치 세 가지 모양입니다. 0이나 마이너스 999처럼 숫자인 척하는 결측치가 가장 위험하고, 이상치는 그림으로 찾되 지울지 남길지는 사람이 판단합니다. 값의 크기가 제각각이면 작은 난수로 시작한 뉴런이 곡선 끝에 붙어 배우지 못합니다. Min-Max나 Z-score로 맞추고, 통계는 학습 데이터로만 구합니다. 무엇마다 맞출지는 데이터가 정합니다. 무엇보다, 어떤 데이터를 받든 모델보다 먼저 데이터를 열어 보세요.}
\end{frame}

% --- 이미지 출처
\begin{frame}{이미지 · 데이터 출처}
\vfill
\small
\begin{itemize}
\item 뉴런 그림: Wikimedia Commons ``Neuron.svg'' (Dhp1080, CC BY-SA 3.0), 라벨 선을 지워 사용
\item 달리는 말: Eadweard Muybridge, \textit{The Horse in Motion} (1878), 퍼블릭 도메인 (Wikimedia Commons)
\item 붓꽃: Iris (Fisher, 1936; scikit-learn 내장) \quad 와인: UCI Wine (Aeberhard \& Forina, 1991) \quad 손글씨 숫자: MNIST (LeCun, Cortes, Burges)
\item 컬러 사진: CIFAR-10, A. Krizhevsky (2009) ``Learning Multiple Layers of Features from Tiny Images''
\item 문자 메시지: UCI SMS Spam Collection (Almeida \& Hidalgo, 2011)
\item 항공 승객: Box \& Jenkins, \textit{Time Series Analysis} (1976), 1949$\sim$1960 국제선 월별 승객
\item 말소리: Free Spoken Digit Dataset (Z. Jackson 외, CC BY-SA 4.0)
\end{itemize}
\vfill
\note{이 덱에 쓴 그림과 데이터의 출처입니다.}
\end{frame}
\end{document}
